\subsection{Spettro dell'oscillatore armonico}
Le relazioni che ho ottenuto sono 
\[
[\hat{a}^\dag, \hat{a}] = -1, \qquad \mathcal{H} = \hbar \omega \left( N + \frac{1}{2} \right).
\]
Voglio determinare lo spettro di $ \mathcal{H} $ e passo attraverso $ [N, \hat{a}] $ e $ [N, \hat{a}^\dag] $.
\[
[N, \hat{a}] = [\hat{a}^\dag \hat{a}, \hat{a}] = [\hat{a}^\dag, \hat{a}]\hat{a} = -\hat{a},
\]
\[
[N, \hat{a}^\dag] = [\hat{a}^\dag \hat{a}, \hat{a}^\dag] = \hat{a}^\dag[\hat{a}, \hat{a}^\dag] = \hat{a}^\dag.
\]

Se conosco un autostato di $ N $, sia $ \ket{z} $, allora
\[
N \ket{z} = \lambda_z \ket{z} \; \Rightarrow \; \begin{cases}
N \hat{a}^\dag \ket{z} = (\lambda_z + 1) \hat{a}^\dag \ket{z} \\
N \hat{a} \ket{z} = (\lambda_z - 1) \hat{a} \ket{z}
\end{cases}.
\]
Dimostrazione:
\begin{align*}
N \hat{a}^\dag \ket{z} &= [N \hat{a}^\dag = \hat{a}^\dag N + \hat{a}^\dag] = (\hat{a}^\dag N + \hat{a}^\dag) \ket{z} = \\
&= \hat{a}^\dag \lambda_z \ket{z} + \hat{a}^\dag \ket{z} = (\lambda_z + 1) \hat{a}^\dag \ket{z} \\
N \hat{a} \ket{z} &= (\hat{a}N - \hat{a}) \ket{z} = (\lambda_z - 1) \hat{a} \ket{z}. \qed
\end{align*}

Inoltre, $ N \ket{z} = \lambda_z \ket{z} \Rightarrow \lambda_z \geq 0 $.

Dimostrazione: $ \braket{z|N|z} = \lambda_z $ quando $ \braket{z|z} = 1 $, quindi
\begin{align*}
\lambda_z &= \braket{z|\hat{a}^\dag \hat{a}|z} = \\
&= \sum_n \braket{z|\hat{a}^\dag|n} \braket{n|\hat{a}|z} = \\
&= \sum_n \left| \braket{n|\hat{a}|z} \right|^2 \geq 0. \qed
\end{align*}

Deve esistere $ \ket{0} $ tale che $ \hat{a} \ket{0} = 0 $, perch\'e se non esistesse, cosiderando
\[
N \hat{a}^P \ket{z} = (\lambda_z - P) \hat{a}^P \ket{z},
\]
per $ P $ abbastanza grande avrei autovalori negativi. Inoltre, se $ \hat{a} \ket{0} = 0 $, allora $ N \ket{0} = \hat{a}^\dag \hat{a} \ket{0} = \lambda_0 \ket{0} $, cio\`e $ \lambda_0 = 0 $.

Allora posso definire
\[
\ket{n} = K_n (\hat{a}^\dag)^n \ket{0}
\]
$ K_n $ costante di normalizzazione tale che 
\[
N \ket{n} = n \ket{n},
\]
\[
\hat{a}^\dag \ket{n} = \frac{K_n}{K_{n + 1}} \ket{n + 1}, \qquad \hat{a} \ket{n} = n \frac{K_n}{K_{n - 1}}\ket{n - 1}.
\]
Dico che $ \ket{0} $ \`e il \textbf{vuoto}, $ N $ \`e l'\textbf{operatore numero}, $ \hat{a} $ e $ \hat{a}^\dag $ sono gli \textbf{operatori di ditruzione e creazione}.

Per calcolare $ K_n $ normalizzo $ \braket{0|0} = 1 $. Allora
\[
\begin{cases}
\ket{n} = K_n (\hat{a}^\dag)^n \ket{0} \\
\braket{n|n} = 1
\end{cases} \quad \Rightarrow
\]
\begin{align*}
\Rightarrow \quad 1 &= |K_n|^2 \braket{0|(\hat{a})^n(\hat{a}^\dag)^n|0} = \\
&= |K_n|^2 \braket{0|(\hat{a})^{n - 1} \hat{a} \hat{a}^\dag (\hat{a}^\dag)^{n - 1}|0} = \\
&= |K_n|^2 \braket{0|(\hat{a})^{n - 1} (\hat{a}^\dag \hat{a} + 1) (\hat{a}^\dag)^{n - 1}|0} = \\
&= |K_n|^2 \braket{0|(\hat{a})^{n - 1} (N + 1) (\hat{a}^\dag)^{n - 1}|0} = \\
&= |K_n|^2 n \braket{0|(\hat{a})^{n - 1} (\hat{a}^\dag)^{n - 1}|0} = \\
&= |K_n|^2 n (n - 1) \braket{0|(\hat{a})^{n - 2} (\hat{a}^\dag)^{n - 2}|0} = \\
&= |K_n|^2 n! \braket{0|0}
\end{align*}
\[
\boxed{K_n = \frac{1}{\sqrt{n!}}} 
\]
da cui segue
\[
\ket{n} = \frac{1}{\sqrt{n!}} (a^\dag)^n \ket{0}
\]
\[
\hat{a}^\dag \ket{n} = \sqrt{n + 1} \ket{n + 1}
\]
\[
\hat{a} \ket{n} = \sqrt{n} \ket{n - 1}
\]
\[
E_n = \hbar \omega \left( n + \frac{1}{2} \right) \text{ autovalori di } \mathcal{H}.
\]

N.B.: il formalismo di Heisenberg ha permesso di riscrivere tutto senza dover scegliere una base.

Gli elementi di matrice li calcolo attraverso $ \hat{a} \pm \hat{a}^\dag $
\[
\hat{x} = \sqrt{\frac{\hbar}{2m\omega}}(\hat{a} + \hat{a}^\dag), \qquad \hat{p} = -i \sqrt{\frac{m\omega\hbar}{2}}(\hat{a} - \hat{a}^\dag)
\]
e usando $ \braket{n|m} = \delta_{nm} $, allora
\[
\begin{cases}
\braket{m|\hat{a}^\dag|n} = \delta_{m, n + 1}\sqrt{n + 1} \\
\braket{m|\hat{a}|n} = \delta_{m, n - 1}\sqrt{n}
\end{cases} \quad \Rightarrow 
\]
\[
\Rightarrow \quad \begin{cases}
\braket{m|\hat{x}|n} = \sqrt{\frac{\hbar}{2m\omega}} (\delta_{m, n - 1}\sqrt{n} + \delta_{m, n + 1}\sqrt{n + 1}) \\
\braket{m|\hat{p}|n} = -i \sqrt{\frac{m\hbar\omega}{2}} (\delta_{m, n - 1}\sqrt{n} - \delta_{m, n + 1}\sqrt{n + 1})
\end{cases}
\]
Allora i valori medi
\[
\braket{n|\hat{x}|n} = \braket{n|\hat{p}|n} = 0.
\]
Per gli operatori al quadrato
\begin{align*}
\braket{n|\hat{x}^2|n} &= \frac{\hbar}{2m\omega} \braket{n|(\hat{a} + \hat{a}^\dag)(\hat{a} + \hat{a}^\dag)|n} = \\
&= \frac{\hbar}{2m\omega} \braket{n|(\hat{a}\hat{a}^\dag + \hat{a}^\dag\hat{a})|n} = \\
&= \frac{\hbar}{2m\omega} \braket{n|(2N + 1)|n} = \\
&= \frac{\hbar}{m\omega} \left( n + \frac{1}{2} \right). \\
\braket{n|\hat{p}^2|n} &= -\frac{m\omega\hbar}{2} \braket{n|(\hat{a}\hat{a}^\dag + \hat{a}^\dag\hat{a})|n} = \\
&= m\omega\hbar \left( n + \frac{1}{2} \right).
\end{align*}

Passo all'energia cinetica e al potenziale (valori medi)
\[
\braket{T} = \frac{\omega\hbar}{2}\left( n + \frac{1}{2} \right),
\]
\[
\braket{V} = \frac{1}{2} m\omega^2 \frac{\hbar}{m\omega} \left( n + \frac{1}{2} \right) = \braket{T},
\]
cio\`e un caso di equipartizione dell'energia.

Le indeterminazioni sono 
\[
\begin{cases}
\braket{n|\Delta^2 \hat{x}|n} = \braket{n|\hat{x}^2|n} \\
\braket{n|\Delta^2 \hat{p}|n} = \braket{n|\hat{p}^2|n}
\end{cases} \quad \Rightarrow
\]
\[
\Rightarrow \quad \Delta^2 \hat{x} \Delta^2 \hat{p} = \hbar^2 \left( n + \frac{1}{2} \right)^2 \geq \frac{\hbar^2}{4}
\]
e $ n = 0 $ d\`a lo stato a minima indeterminazione.